\section{From Patrol Coverage to Loss-Flow Interdiction}
\subsection{Why Coverage Maps Miss the Hard Part of Protection}
Large protected areas are often discussed as if protection quality were proportional to the fraction of land that receives some visible patrol presence. That logic is too coarse for real operations. Illegal actors do not value every location equally, do not enter the park uniformly, and do not experience the same response difficulty everywhere. In practice, ecological damage is transmitted through a sparse set of pathways: entry opportunities, movement corridors, relay zones, and high-value habitat clusters. If those pathways remain weakly monitored, a park can look busy on paper while still leaking high-consequence losses in the field \citep{watson2014,geldmann2013,coad2019,jones2018}.

Etosha National Park is a strong example of this mismatch. Wildlife value is concentrated around water-dependent habitat, road access is uneven, and station geometry creates service asymmetry across the landscape \citep{etoshaOfficial,etoshaEcology}. Those ingredients imply that the key management task is not to color the map with generalized effort, but to suppress a small set of biologically expensive pathways before they mature into realized loss.

\subsection{Problem Restatement}
We therefore restate the IMMC problem as follows: build a decision system that minimizes residual weighted wildlife-loss flow from ingress points to high-value habitat under limited patrol teams, limited sensing infrastructure, finite drone support, and uncertain daily pressure.

This is a materially different control problem from standard hotspot ranking. The manager is not merely selecting where risk is high; the manager is choosing which pathways will remain open after a protection plan is deployed. Let
\[
\mathcal{N}_t = (\mathcal{E},\mathcal{Z},\mathcal{H},\mathcal{R},\Lambda_t)
\]
denote the park network at time \(t\), where \(\mathcal{E}\) is the ingress set, \(\mathcal{Z}\) the management zones, \(\mathcal{H}\) the habitat sinks, \(\mathcal{R}\) the patrol road system, and \(\Lambda_t\) the current pressure and service state. The control decision
\[
u_t=(x_t,y_t,z_t)
\]
assigns patrol templates, sentinel devices, and UAV reserve windows. The objective is to minimize the residual weighted loss-flow that survives those controls while keeping service delays within a defensible commitment standard.

\subsection{New Analytic Spine}
To operationalize that restatement, we use a three-core architecture plus one communication layer:
\begin{enumerate}
  \item a \textbf{threat-transport graph} that estimates how pressure enters the park and how it propagates toward high-value habitat;
  \item a \textbf{lexicographic barrier-design model} that caps the worst remaining hotspot debt before minimizing total residual loss;
  \item a \textbf{queueing-based reliability model} that translates station load and travel time into a staffing requirement;
  \item a \textbf{managerial synthesis layer} that converts the mathematics into an operational doctrine.
\end{enumerate}

The information flow is
\[
\text{basemap + roads + stations + habitat}
\;\longrightarrow\;
\text{transport graph + hinge zones}
\;\longrightarrow\;
\text{barrier design}
\;\longrightarrow\;
\text{service reliability and staffing}.
\]
Figure~\ref{fig:system} summarizes this pipeline visually. The central unit of decision is not a generic hotspot, but a pathway bottleneck with measurable residual debt.

\begingroup
\setstretch{0.6}
\begin{figure}[H]
  \centering
  \includegraphics[width=0.8\linewidth]{fig01_system_architecture.png}
  \caption{Loss-flow interdiction architecture linking transport-graph diagnosis, barrier design, service-reliability sizing, and management feedback.}
  \label{fig:system}
\end{figure}
\endgroup

\subsection{Management Metrics}
We evaluate plans with three cross-functional metrics:
\begin{align}
\PI &: \text{risk-weighted protection index, i.e., the complement of residual weighted loss flow}, \\
\Gamma_{\max} &: \text{worst residual debt among hinge zones}, \\
\Pr(T>\tau^\star) &: \text{delay exceedance probability relative to the response commitment}.
\end{align}
These indicators were chosen to force the analysis to stay operational. A plan is not attractive if it looks efficient on average but leaves one hinge zone catastrophically exposed or pushes one station into chronic overload.

\subsection{What This Changes for the IMMC Tasks}
The seven IMMC tasks are still answered, but under a different logic than a standard patrol-allocation paper. Table~\ref{tab:qmap} shows the mapping.
\begin{table}[H]
  \centering
  \caption{Mapping from IMMC tasks to the re-derived analytic outputs.}
  \label{tab:qmap}
  \begin{tabular}{p{0.12\linewidth}p{0.27\linewidth}p{0.45\linewidth}}
    \toprule
    Task & Decision target & Output in this paper \\
    \midrule
    Task 1 & Define protection and identify priorities & Weighted loss-flow field, hinge-zone diagnosis, and response commitment \\
    Task 2 & Allocate resources & Lexicographic barrier design across patrol templates, sentinels, and UAV reserve \\
    Task 3 & Predict performance over time and staffing & Queueing reliability envelope plus backtested protection index \\
    Task 4 & Sensitivity & Structural trade-off frontier in manpower, uptime, and cadence \\
    Task 5 & Assess model quality & Identification risks, adversarial failure cases, and validation boundary \\
    Task 6 & Adapt to other parks & Morphology-based transfer protocol using ingress, sink, and service-rate rebuilding \\
    Task 7 & Communicate results & Letter and summary sheet framed in management language rather than solver language \\
    \bottomrule
  \end{tabular}
\end{table}
